% !TEX encoding = UTF-8 Unicode
% ---------------------------------------------------------------------
% 7th Photomechanics Conference, 5th IDICs Conference
% ---------------------------------------------------------------------
% 19-22 october 2020, Nantes, France
% ---------------------------------------------------------------------
\documentclass{PMIDIC2020}

% ---------------------------------------------------------------------
\title{\LaTeX ~Template for PM-IDICs 2021 abstracts}
% ---------------------------------------------------------------------
\author{M. Coret, M.L.M. Fran\c cois, J. R\'ethor\'e, R. Seghir}
% ---------------------------------------------------------------------
\address{%
Centrale Nantes, Universit\'e de Nantes, GeM (UMR6183), Nantes, France
}
% ---------------------------------------------------------------------
\begin{document}
% ---------------------------------------------------------------------
\maketitle
% ---------------------------------------------------------------------

\begin{abstract}
Here you write your 5 to 6 lines abstract that summarizes the \underline{2 pages article}.\\
\vspace{1.5cm}\\
\keywords DIC, PhotoMechanics
\end{abstract}



\paragraph{Introduction}
As much as possible, we encourage you to write your article using the IMRAD method. 
\begin{center}
\url{https://en.wikipedia.org/wiki/IMRAD}
\end{center}
\vspace{1.5cm}

\paragraph{Methods}
Explanations of the method ...\\

If you have equations, they are centered and numbered. For example see  eq.\eqref{eq:bc}.
\begin{equation}
\mathbf u(\mathbf x)=\arg\min \int \left[ f(\mathbf x)-g(\mathbf x+\mathbf u(\mathbf x))\right]^2 d\mathbf x \label{eq:bc}
\end{equation}

\vspace{1.5cm}

\paragraph{Results}
Development of the results ...\\

If you have figures, they are centered and numbered. For example see fig.\eqref{fig:pfe}.

\begin{figure}[htbp!]
\centering
\includegraphics[width=0.3\linewidth]{disp_field_example}
\caption{Horizontal displacement field of ...}\label{fig:pfe}
\end{figure}

\paragraph{Discussion and Conclusion} This is the place to discuss the results and give conclusions and perspectives ...\\

If you have tabulars, they are centered and numbered. For example see tab.\eqref{tab:nice}.

\begin{table}[htbp!]
\centering
\begin{tabular}{|c|c|c|}
  \hline
\textbf{  column 1} & \textbf{column 2 }&\textbf{ column 3 }\\
  \hline
  1.1 & 1.2 & 1.3 \\
  2.1 & 2.2 & 2.3 \\
  \hline
\end{tabular}
\caption{Table of ... }\label{tab:nice}
\end{table}


%

References and citations have to be written like this: 
\begin{description}
\item[$\bullet$] citation of a journal paper \cite{horn81},
\item[$\bullet$] citation of a conference proceeding \cite{lucas81}, 
\item[$\bullet$] citation of  a book \cite{sutton09}.
\end{description}


% ---------------------------------------------------------------------
%
% ---------------------------------------------------------------------
%\\ article
\begin{thebibliography}{1}

\bibitem{horn81}
B.K.P. Horn and G.~Schunck.
\newblock Determining optical flow.
\newblock {\em Artificial Intelligence}, 17:185--203, 1981.

\bibitem{lucas81}
B.D. Lucas and T.~Kanade. 
\newblock An iterative image registration technique with an application to stereo vision. 
\newblock In {\em Proceedings of Imaging Understanding Workshop}, pages 121--130, 1981.

\bibitem{sutton09}
M.A. Sutton, J.-J.~Orteu, H.~Schreier, 
\newblock \emph{Image correlation for shape, motion and deformation measurements: Basic Concepts, Theory and Applications}, 
\newblock Springer, New York, NY (USA), 2009.


% ---------------------------------------------------------------------
\end{thebibliography}
% ---------------------------------------------------------------------

% ---------------------------------------------------------------------
\end{document}
% ---------------------------------------------------------------------


